% Chapitre 6 : écran fictif de la page de relecture locale.
\begin{figure}[tbp]
\centering
\begin{tikzpicture}[x=1cm,y=1cm,
  texte/.style={font=\sffamily\scriptsize, anchor=west, inner sep=0pt}]
  \fenetre{12.4}{8.0}{Relecture}
  % barre d'adresse
  \draw[black!25, fill=black!3, rounded corners=5pt] (0.3,6.95) rectangle (12.1,7.3);
  \node[font=\ttfamily\scriptsize, text=black!60, anchor=west] at (0.5,7.125)
    {localhost:5173/relecture};
  % le fac-similé
  \draw[black!40, pattern=north east lines, pattern color=black!18] (0.3,0.4) rectangle (5.6,6.6);
  \node[fill=white, font=\sffamily\scriptsize, text=black!60, inner sep=3pt] at (2.95,4.9)
    {page scannée};
  \draw[tonregle, line width=1pt, fill=white, fill opacity=0.6] (1.5,3.0) rectangle (3.9,3.5);
  \node[font=\sffamily\tiny, text=tonregle, anchor=north west] at (1.5,2.95) {le mot en question};
  \node[font=\sffamily\tiny, text=black!55, anchor=south west] at (0.35,0.45)
    {lot 4, page 7};
  % la carte de question
  \draw[black!35, rounded corners=3pt, fill=white] (6.0,1.6) rectangle (12.1,6.6);
  \node[texte, text=black!55] at (6.25,6.25) {Doute 3 sur 12};
  \fill[black!10] (9.6,6.2) rectangle (11.85,6.3);
  \fill[tonregle] (9.6,6.2) rectangle (10.15,6.3);
  \node[texte, font=\sffamily\scriptsize\bfseries] at (6.25,5.7) {Lecture proposée};
  \node[texte, font=\ttfamily\scriptsize] at (6.25,5.3) {\textbackslash uncertain\{faisceau cohérent\}};
  \node[texte, anchor=north west, text width=5.6cm, align=left] at (6.25,4.85)
    {Le mot encadré sur la page se lit-il bien \mbox{\og cohérent\fg{} ?}};
  \node[bouton, anchor=west] (oui) at (6.25,3.85) {Oui};
  \node[bouton, anchor=west, right=2mm of oui] (non) {Non};
  \node[bouton, anchor=west, right=2mm of non] {Incertain};
  \draw[black!30, rounded corners=2pt] (6.25,2.0) rectangle (11.85,3.2);
  \node[texte, text=black!40, anchor=north west] at (6.4,3.05) {Note (facultative)};
  % pied de page
  \node[texte, text=black!55] at (6.05,1.05) {9 réponses en attente};
  \node[bouton, fill=tonregle, text=white, draw=tonregle, anchor=east] at (12.1,1.05)
    {Copier les réponses};
\end{tikzpicture}
\caption{Écran fictif : la page de relecture locale. À gauche la page de l'original, à
droite une question fermée sur la lecture proposée. Les réponses s'accumulent dans le
navigateur jusqu'à ce qu'on les copie dans la conversation.}
\label{fig:ecran-relecture}
\end{figure}
